\input{../../../templates/es/preambulo.tex}

% Los listados se toman de src/ con \lstinputlisting; la diapositiva es el archivo.
\lstset{
    basicstyle=\ttfamily\fontsize{8pt}{9.6pt}\selectfont,
    breaklines=true,
    breakatwhitespace=true,
    columns=fullflexible,
    keepspaces=true,
    xleftmargin=0.3em,
    framesep=2pt,
    aboveskip=0pt,
    belowskip=0pt
}

\newcommand{\marcoresultado}[3]{%
    \begin{frame}{#1}
        \begin{center}
            \includegraphics[height=0.62\textheight,width=\textwidth,keepaspectratio]{#2}
        \end{center}
        \vspace{0.25em}
        {\small #3\par}
    \end{frame}%
}

\newcommand{\marcoafirmacion}[2]{%
    \begin{frame}{#1}
        {\small #2\par}
    \end{frame}%
}

\date{}

\begin{document}

\tituloleccion{10}{Optimización Combinatoria Ordenada}

\section{Esta lección}

\begin{frame}{Esta lección}
La Lección 09 eran bits. Un tour no es una máscara: la cruza de un punto duplica ciudades. La cruza ordenada es el operador que sigue siendo legal.

\vspace{0.6em}
\begin{itemize}
    \item Un cromosoma ordenado legal contiene cada miembro requerido exactamente una vez.
    \item La cruza ordinaria de un punto suele duplicar y perder miembros.
    \item Un resultado válido de AG todavía necesita una línea base (vecino más cercano; un presupuesto).
\end{itemize}
\end{frame}

% ================= equipo_01_la_plantilla =================
\begin{frame}{Dos modelos de permutación restringida}
Para una ruta \(\pi\), con \(\pi_{n+1}=\pi_1\),
\[
L(\pi)=\sum_{i=1}^{n}d(\pi_i,\pi_{i+1}).
\]
OX1 e inversión conservan pertenencia y unicidad. Una plantilla legal usa
identificadores distintos en posiciones por rol, respeta el presupuesto y
maximiza habilidad. Estos resultados certifican legalidad y comparan líneas
base; no prueban optimalidad global.
\end{frame}

\section{Once espacios ordenados}

\begin{frame}{Once espacios ordenados}
El cromosoma es una lista ordenada de once identificadores distintos de jugadores. El orden
de los espacios codifica un portero, cuatro defensas, tres mediocampistas y tres
delanteros, así que tanto la pertenencia como la posición deben sobrevivir a todo operador.
\end{frame}

\marcoafirmacion{Once espacios ordenados}{%
    El grupo respaldado por la fuente contiene 12 candidatos para cada bloque de rol y un presupuesto explícito de €600m}

% ================= equipo_02_plantillas_aleatorias =================
\section{Roles legales aleatorios}

\begin{frame}{Roles legales aleatorios}
Roles legales aleatorios
\end{frame}

\begin{frame}[fragile]{(1) equipo\_aleatorio()}
\lstinputlisting[basicstyle=\ttfamily\fontsize{8pt}{9.6pt}\selectfont,firstline=47,lastline=54]{../src/equipo_02_plantillas_aleatorias.py}
\end{frame}

\begin{frame}[fragile]{(2) metricas\_equipo()}
\lstinputlisting[basicstyle=\ttfamily\fontsize{8pt}{9.6pt}\selectfont,firstline=57,lastline=61]{../src/equipo_02_plantillas_aleatorias.py}
\end{frame}

\begin{frame}[fragile]{(3) el censo}
\lstinputlisting[basicstyle=\ttfamily\fontsize{8pt}{9.6pt}\selectfont,firstline=64,lastline=70]{../src/equipo_02_plantillas_aleatorias.py}
\end{frame}

\marcoafirmacion{Roles legales aleatorios}{%
    Todas las 1,000 plantillas contienen once jugadores distintos; solo 440 cumplen con el presupuesto}

% ================= equipo_03_reparacion_presupuesto =================
\section{Reparación de presupuesto}

\begin{frame}{Reparación de presupuesto}
Reparación de presupuesto
\end{frame}

\begin{frame}[fragile]{(1) reparar\_presupuesto()}
\lstinputlisting[basicstyle=\ttfamily\fontsize{8pt}{9.6pt}\selectfont,firstline=58,lastline=75]{../src/equipo_03_reparacion_presupuesto.py}
\end{frame}

\begin{frame}[fragile]{(2) la auditoria}
\lstinputlisting[basicstyle=\ttfamily\fontsize{8pt}{9.6pt}\selectfont,firstline=78,lastline=85]{../src/equipo_03_reparacion_presupuesto.py}
\end{frame}

\marcoafirmacion{Reparación de presupuesto}{%
    La reparación hace que 1,000/1,000 equipos sean asequibles y preserva la unicidad}

% ================= equipo_04_operadores_ordenados =================
\section{Operadores ordenados}

\begin{frame}{Operadores ordenados}
Operadores ordenados
\end{frame}

\begin{frame}[fragile]{(1) cruza\_roles()}
\lstinputlisting[basicstyle=\ttfamily\fontsize{8pt}{9.6pt}\selectfont,firstline=70,lastline=78]{../src/equipo_04_operadores_ordenados.py}
\end{frame}

\begin{frame}[fragile]{(2) mutar()}
\lstinputlisting[basicstyle=\ttfamily\fontsize{8pt}{9.6pt}\selectfont,firstline=81,lastline=86]{../src/equipo_04_operadores_ordenados.py}
\end{frame}

\begin{frame}[fragile]{(3) una generacion}
\lstinputlisting[basicstyle=\ttfamily\fontsize{8pt}{9.6pt}\selectfont,firstline=89,lastline=94]{../src/equipo_04_operadores_ordenados.py}
\end{frame}

\marcoafirmacion{Operadores ordenados}{%
    Todos los 100 hijos siguen siendo legales; la habilidad media sube de 979.6 a 985.3}

% ================= equipo_05_la_busqueda_completa =================
\section{La búsqueda completa de plantilla}

\begin{frame}{La búsqueda completa de plantilla}
La búsqueda completa de plantilla
\end{frame}

\begin{frame}[fragile]{(1) ejecutar()}
\lstinputlisting[basicstyle=\ttfamily\fontsize{8pt}{9.6pt}\selectfont,firstline=99,lastline=112]{../src/equipo_05_la_busqueda_completa.py}
\end{frame}

\begin{frame}[fragile]{(2) la linea base}
\lstinputlisting[basicstyle=\ttfamily\fontsize{8pt}{9.6pt}\selectfont,firstline=115,lastline=118]{../src/equipo_05_la_busqueda_completa.py}
\end{frame}

\begin{frame}[fragile]{(3) la plantilla}
\lstinputlisting[basicstyle=\ttfamily\fontsize{8pt}{9.6pt}\selectfont,firstline=121,lastline=124]{../src/equipo_05_la_busqueda_completa.py}
\end{frame}

\marcoresultado{La búsqueda completa de plantilla}{../figures/equipo_05_la_busqueda_completa.png}{%
    El AG alcanza una habilidad de 991 con €597.5m y empata a la mejor plantilla aleatoria reparada a iguales evaluaciones}

% ================= tsp_01_la_ruta =================
\section{Una ruta es un ordenamiento}

\begin{frame}{Una ruta es un ordenamiento}
Un cromosoma del TSP contiene cada ciudad exactamente una vez. Sus genes no son opciones separadas: cambiar una posición cambia el significado de toda la ruta.
\end{frame}

\marcoafirmacion{Una ruta es un ordenamiento}{%
    La instancia proporcionada contiene 48 ciudades; la ruta identidad es legal pero larga}

% ================= tsp_02_poblacion_aleatoria =================
\section{Rutas legales aleatorias}

\begin{frame}{Rutas legales aleatorias}
Rutas legales aleatorias
\end{frame}

\begin{frame}[fragile]{(1) ruta\_aleatoria()}
\lstinputlisting[basicstyle=\ttfamily\fontsize{8pt}{9.6pt}\selectfont,firstline=35,lastline=40]{../src/tsp_02_poblacion_aleatoria.py}
\end{frame}

\begin{frame}[fragile]{(2) TAMANO\_POBLACION}
\lstinputlisting[basicstyle=\ttfamily\fontsize{8pt}{9.6pt}\selectfont,firstline=43,lastline=45]{../src/tsp_02_poblacion_aleatoria.py}
\end{frame}

\begin{frame}[fragile]{(3) el censo}
\lstinputlisting[basicstyle=\ttfamily\fontsize{8pt}{9.6pt}\selectfont,firstline=48,lastline=53]{../src/tsp_02_poblacion_aleatoria.py}
\end{frame}

\marcoafirmacion{Rutas legales aleatorias}{%
    Todas las 500 rutas barajadas son legales, con una amplia distribución de longitudes}

% ================= tsp_03_cruza_ordenada =================
\section{La cruza debe preservar una permutación}

\begin{frame}{La cruza debe preservar una permutación}
La cruza debe preservar una permutación
\end{frame}

\begin{frame}[fragile]{(1) un\_punto()}
\lstinputlisting[basicstyle=\ttfamily\fontsize{8pt}{9.6pt}\selectfont,firstline=39,lastline=43]{../src/tsp_03_cruza_ordenada.py}
\end{frame}

\begin{frame}[fragile]{(2) cruza\_ordenada()}
\lstinputlisting[basicstyle=\ttfamily\fontsize{8pt}{9.6pt}\selectfont,firstline=46,lastline=56]{../src/tsp_03_cruza_ordenada.py}
\end{frame}

\begin{frame}[fragile]{(3) el experimento}
\lstinputlisting[basicstyle=\ttfamily\fontsize{8pt}{9.6pt}\selectfont,firstline=59,lastline=68]{../src/tsp_03_cruza_ordenada.py}
\end{frame}

\marcoafirmacion{La cruza debe preservar una permutación}{%
    La cruza de un punto produce 0/1,000 hijos legales; la cruza ordenada produce 1,000/1,000}

% ================= tsp_04_mutacion =================
\section{La mutación cambia el orden, no la pertenencia}

\begin{frame}{La mutación cambia el orden, no la pertenencia}
La mutación cambia el orden, no la pertenencia
\end{frame}

\begin{frame}[fragile]{(1) mutacion\_intercambio()}
\lstinputlisting[basicstyle=\ttfamily\fontsize{8pt}{9.6pt}\selectfont,firstline=48,lastline=54]{../src/tsp_04_mutacion.py}
\end{frame}

\begin{frame}[fragile]{(2) mutacion\_inversion()}
\lstinputlisting[basicstyle=\ttfamily\fontsize{8pt}{9.6pt}\selectfont,firstline=57,lastline=63]{../src/tsp_04_mutacion.py}
\end{frame}

\begin{frame}[fragile]{(3) la auditoria de distancia}
\lstinputlisting[basicstyle=\ttfamily\fontsize{8pt}{9.6pt}\selectfont,firstline=66,lastline=73]{../src/tsp_04_mutacion.py}
\end{frame}

\marcoafirmacion{La mutación cambia el orden, no la pertenencia}{%
    El intercambio y la inversión preservan cada ciudad; sus cambios medios de longitud son +451.0 y +238.9}

% ================= tsp_05_la_busqueda_completa =================
\section{La búsqueda ordenada completa}

\begin{frame}{La búsqueda ordenada completa}
La búsqueda ordenada completa
\end{frame}

\begin{frame}[fragile]{(1) vecino\_mas\_cercano()}
\lstinputlisting[basicstyle=\ttfamily\fontsize{8pt}{9.6pt}\selectfont,firstline=57,lastline=65]{../src/tsp_05_la_busqueda_completa.py}
\end{frame}

\begin{frame}[fragile]{(2) seleccion\_torneo()}
\lstinputlisting[basicstyle=\ttfamily\fontsize{8pt}{9.6pt}\selectfont,firstline=68,lastline=71]{../src/tsp_05_la_busqueda_completa.py}
\end{frame}

\begin{frame}[fragile]{(3) ejecutar()}
\lstinputlisting[basicstyle=\ttfamily\fontsize{8pt}{9.6pt}\selectfont,firstline=74,lastline=91]{../src/tsp_05_la_busqueda_completa.py}
\end{frame}

\marcoresultado{La búsqueda ordenada completa}{../figures/tsp_05_la_busqueda_completa.png}{%
    La ruta del AG es legal pero 57.2\% más larga que el vecino más cercano después de 10,100 evaluaciones}

\section{Siguiente}

\begin{frame}{Siguiente}
Esta instancia del TSP es la que reutilizará la Lección 12, reescrita en local. Es la única importación genuina entre capítulos del libro.

\vspace{0.8em}
\textbf{Lección 11 --- Otros problemas}

Aquí ambos resultados se comprueban por legalidad y contra líneas base nombradas. Ninguna comparación prueba un óptimo global --- justo la distinción que exige la Lección 08.
\end{frame}
\end{document}
